\documentclass[11pt,twocolumn]{article}
\usepackage[margin=0.85in,top=1in]{geometry}

\usepackage{amsmath,amssymb,amsthm}
\usepackage{graphicx}
\usepackage{microtype}
\usepackage{hyperref}
\usepackage{xcolor}
\usepackage{booktabs}
\usepackage{titlesec}
\titleformat{\section}{\normalfont\bfseries\large}{\thesection}{1em}{}
\titleformat{\subsection}{\normalfont\bfseries\normalsize}{\thesubsection}{1em}{}
\titlespacing*{\section}{0pt}{14pt}{4pt}
\titlespacing*{\subsection}{0pt}{10pt}{3pt}
\hypersetup{colorlinks=true,linkcolor=blue,citecolor=blue,urlcolor=blue}

% ── Shortcuts ──────────────────────────────────────────────────────────
\newcommand{\ket}[1]{|#1\rangle}
\newcommand{\bra}[1]{\langle #1|}
\newcommand{\braket}[2]{\langle #1|#2\rangle}
\newcommand{\rhoS}{\rho_S}
\newcommand{\DS}{D_{\mathrm{S}}}
\newcommand{\Drel}{D(\rho\,\|\,\sigma_0)}
\newcommand{\DrelS}{S(\rho\,\|\,\sigma_0)}
\newcommand{\tauact}{\tau_{\mathrm{act}}}
\newcommand{\vLR}{v_{\mathrm{LR}}}
\newcommand{\PDG}{\textsc{pdg}}

\begin{document}

\title{Wave Function Collapse as Irreversible Entropy Production:\\
An Observer-Independent Criterion with Three Experimental Predictions}

\author{Joshua F.~Sandeman\\
\small Independent Researcher, Salem, Oregon, USA\\
\small \href{mailto:sansuikyo@gmail.com}{sansuikyo@gmail.com}}

\date{March 2026}

% ── ABSTRACT ───────────────────────────────────────────────────────────
\begin{abstract}
The quantum measurement problem asks why a definite outcome occurs
when quantum mechanics predicts only a probability distribution.
We present an observer-independent criterion for wave function
collapse: an interaction constitutes a genuine actualization event
if and only if it produces an irreversible increase in quantum relative
entropy $\Delta S(\rho\|\sigma_0) > 0$ with respect to the Minkowski
vacuum state $\sigma_0$.
This criterion is basis-independent, non-perturbative, and does
not require an observer, a Heisenberg cut, or a preferred foliation.
In weakly coupled quantum field theory it coincides with on-shell
emission ($p^\mu p_\mu = m^2 c^2$), recovering the standard perturbative
picture as a limiting case.
The framework resolves the measurement problem by dissolving it:
the question ``why does a definite outcome occur?'' is answered by
``because an irreversible entropy-producing event occurred,'' which
is a physical fact about the interaction, not a postulate about
observers.
We derive three experimental predictions that discriminate
this framework from all current interpretations of quantum mechanics:
(1) a finite actualization propagation velocity bounded by the
Lieb-Robinson speed; (2) an environment-dependent transition
duration scaling as $\tauact \sim \hbar/\Delta E_{\mathrm{env}}$;
and (3) a sharp, measurable reversibility threshold above which
quantum transitions cannot be undone regardless of intervention.
All three are absent from dynamical collapse models, decoherence-only
accounts, and Everettian approaches.
The Rindler thermal bath satisfies $\Delta S(\rho\|\sigma_0) = 0$
under modular flow, showing that a uniformly accelerating observer
detects no actualization events in the Minkowski vacuum---a result
formally verified in Lean~4/Mathlib.
\end{abstract}

\maketitle

% ── I. INTRODUCTION ────────────────────────────────────────────────────
\section{Introduction}

The measurement problem is the most persistent foundational puzzle
in quantum mechanics~\cite{Bell1990}.
The Schr\"odinger equation is unitary and linear: it evolves
superpositions into larger superpositions without selecting an
outcome. Yet every experiment produces a definite outcome.
Reconciling these two facts requires either a modification of the
dynamics (dynamical collapse~\cite{GRW1986,Pearle1989}), a reinterpretation
of outcomes (Everett~\cite{Everett1957}, relational QM~\cite{Rovelli1996}),
or an additional ontological element (pilot wave~\cite{Bohm1952}).

None of these approaches offers an observer-independent physical
criterion for \emph{when} collapse occurs.
Dynamical collapse models introduce new free parameters
(the GRW collapse rate $\lambda$, the correlation length $r_C$)
and have been constrained but not verified by experiment.
Everettian approaches deny that collapse occurs at all.
Relational QM grounds collapse in observations but does not
specify what constitutes an observation.
Decoherence accounts correctly identify environmental entanglement
as the mechanism for apparent classicality but, as Joos and
Zeh~\cite{Joos1985} and later Zurek~\cite{Zurek2003} both acknowledged,
decoherence does not by itself select a definite outcome.

We propose a criterion that is physical rather than observer-relative,
non-perturbative rather than tied to Feynman diagrams, and
falsifiable rather than merely interpretive.
The central claim is stated in Section~\ref{sec:criterion}.
Sections~\ref{sec:resolution}--\ref{sec:predictions} develop
its consequences: the dissolution of the measurement problem,
the Unruh consistency check, and three discriminating experimental
predictions.
Section~\ref{sec:discussion} positions the framework against
competing approaches.

% ── II. THE ACTUALIZATION CRITERION ────────────────────────────────────
\section{The Actualization Criterion}
\label{sec:criterion}

\subsection{Setup}

Let $\mathcal{H}$ be the Hilbert space of a quantum system $S$
coupled to an environment $E$. The composite state evolves under
the Schr\"odinger equation:
\begin{equation}
  \ket{\Psi(t)} = \alpha\ket{x_1}\ket{\varepsilon_1(t)}
                + \beta\ket{x_2}\ket{\varepsilon_2(t)},
  \label{eq:composite}
\end{equation}
where $\ket{\varepsilon_1}$, $\ket{\varepsilon_2}$ are environmental
states correlated with the system branches.
The reduced density matrix of $S$ is
$\rhoS(t) = \operatorname{Tr}_E\bigl[\ket{\Psi(t)}\bra{\Psi(t)}\bigr]$.
The coherences are suppressed by the overlap
$\braket{\varepsilon_2(t)}{\varepsilon_1(t)} \to 0$
as the environment records which-branch information.

Let $\sigma_0$ denote the Minkowski vacuum state of the field.
The quantum relative entropy of $\rhoS$ with respect to $\sigma_0$
is~\cite{Umegaki1962,Lindblad1975}
\begin{equation}
  S(\rhoS\|\sigma_0)
  = \operatorname{Tr}\bigl[\rhoS(\ln \rhoS - \ln \sigma_0)\bigr].
  \label{eq:relent}
\end{equation}

\subsection{The criterion}
\label{sec:criterion_statement}

\begin{quote}
\textbf{Actualization Criterion.}
An interaction between system $S$ and environment $E$ constitutes
a genuine \emph{actualization event}---a wave function collapse in
the ontologically real sense---if and only if it produces an
irreversible increase
\begin{equation}
  \Delta S(\rhoS\|\sigma_0) > 0
  \label{eq:criterion}
\end{equation}
in the quantum relative entropy of $S$ with respect to the
vacuum state $\sigma_0$.
\end{quote}

Several features of this criterion are immediate.

\emph{Basis-independence.}
The relative entropy~\eqref{eq:relent} is defined at the level
of the density matrix and does not depend on the choice of
computational basis.

\emph{Observer-independence.}
The criterion refers to a physical property of the interaction
($\Delta S > 0$), not to any observer's record or knowledge.
The vacuum state $\sigma_0$ is a physical state of the field,
not a subjective reference.

\emph{Non-perturbative.}
The criterion is formulated in terms of the density matrix and
applies regardless of coupling strength, extending beyond the
domain of Feynman perturbation theory.

\emph{Irreversibility as the key.}
The condition is $\Delta S > 0$, not merely $\Delta S \ne 0$.
A reversible interaction (such as a coherent scattering event
that does not leave a lasting environmental record) satisfies
$\Delta S = 0$ and is not an actualization event. This
distinguishes the criterion from entropy production in
non-equilibrium thermodynamics: not all entropy-producing
processes constitute actualizations, but all actualizations
are entropy-producing.

\subsection{On-shell condition as perturbative limit}

In weakly coupled quantum field theory, an interaction in which
a real, on-shell excitation satisfying $p^\mu p_\mu = m^2 c^2$
propagates into the asymptotic environment is a \emph{sufficient}
condition for~\eqref{eq:criterion}.
An on-shell emission creates environmental states that are
asymptotically orthogonal and generates a lasting, redundant
causal record.
Purely virtual (off-shell) exchange satisfies $\Delta S = 0$
and is not an actualization.
The on-shell condition is therefore the perturbative implementation
of the entropic criterion, not a separate postulate.

% ── III. DISSOLUTION OF THE MEASUREMENT PROBLEM ────────────────────────
\section{Resolution of the Measurement Problem}
\label{sec:resolution}

\subsection{The puzzle dissolved}

The measurement problem asks: why does the system take on a
definite value when quantum mechanics predicts only a probability
distribution?

The answer within this framework is: because an actualization
event satisfying~\eqref{eq:criterion} has occurred.
The transition from superposition to definite outcome is not a
postulate; it is the physical fact that the interaction produced
an irreversible increase in relative entropy.
The question is dissolved rather than solved: asking
``why did collapse occur?'' is equivalent to asking ``why did this
interaction produce a permanent environmental record?''---which is
a question about interaction physics, not about the foundations
of quantum mechanics.

\subsection{Schr\"odinger's cat}

The cat thought experiment is often presented as showing that
quantum superpositions can persist at macroscopic scales.
Within this framework, a real cat is continuously actualizing:
it metabolizes, moves, and emits heat---a dense internal network
of on-shell interactions that pins it firmly in the classical regime.
The genuine locus of quantum indeterminacy is the isotope: a small,
isolated nucleus whose decay satisfies~\eqref{eq:criterion} at
the moment of emission.
The cat is never in superposition because it is a densely coupled
actualization network, not a microscopic superposition.
The paradox dissolves.

\subsection{The Heisenberg cut is unnecessary}

Decoherence-based accounts of measurement require a Heisenberg
cut---a boundary at which quantum mechanics hands off to classical
physics. The criterion~\eqref{eq:criterion} removes the need for
such a cut.
The classical regime is the dense-graph limit of the same dynamics:
as the environmental coupling strength increases,
$\braket{\varepsilon_2(t)}{\varepsilon_1(t)} \to 0$ rapidly,
and~\eqref{eq:criterion} is satisfied on timescales far shorter
than any observable interval.
There is no boundary between quantum and classical regimes,
only a gradient in the local actualization rate.

\subsection{Unruh consistency}

A uniformly accelerating observer (Rindler observer) perceives
the Minkowski vacuum as a thermal bath at temperature
$T_U = \hbar a / (2\pi c k_B)$~\cite{Unruh1976}.
A potential paradox arises: does the Rindler thermal bath trigger
actualization events in the Minkowski vacuum?

The criterion~\eqref{eq:criterion} resolves this cleanly.
The Rindler thermal bath is the Minkowski vacuum expressed in
the Rindler algebra; under modular flow, the relative entropy
satisfies $\Delta S(\rho\|\sigma_0) = 0$~\cite{Haag1996}.
No actualization event occurs.
Both the inertial observer (who sees $|0_M\rangle$) and the
Rindler observer (who sees a thermal bath) agree: there are no
physical snap events in the Minkowski vacuum.
This result has been formally verified in Lean~4/Mathlib as the
theorem \texttt{rindler\_stationarity} in
\texttt{RA\_AQFT\_Proofs\_v10.lean}:
the Rindler density matrix $\rho_\beta$ satisfies
$S(\rho_\beta\|\sigma_0) = 0$ under modular flow.

% ── IV. THREE EXPERIMENTAL PREDICTIONS ─────────────────────────────────
\section{Three Experimental Predictions}
\label{sec:predictions}

The criterion~\eqref{eq:criterion} makes three quantitative
predictions that are absent from all current competing frameworks.
Table~\ref{tab:predictions} summarizes the discriminating power
against four alternative approaches.

\subsection{Prediction 1: Finite actualization propagation velocity}

The criterion~\eqref{eq:criterion} requires a physical mechanism
to propagate from the locus of the interaction to all regions
whose superposition is terminated.
This mechanism is the spread of environmental entanglement,
which is bounded by the Lieb-Robinson speed~\cite{LiebRobinson1972}.

\textbf{Prediction.}
In a spatially extended quantum many-body system with
nearest-neighbor coupling $J$ and lattice spacing $\ell$, the
decoherence of a distant site $B$ following an actualization event
at site $A$ exhibits a time delay
\begin{equation}
  \Delta t \geq \frac{|r_A - r_B|}{\vLR},
  \qquad \vLR \sim \frac{J\ell}{\hbar},
  \label{eq:LR}
\end{equation}
before the coherences at $B$ begin to collapse.
Instantaneous collapse models (GRW, CSL) predict $\Delta t = 0$.
Everettian and decoherence-only accounts predict that
``collapse'' at $B$ is merely apparent and has no physical
propagation velocity.
The prediction~\eqref{eq:LR} is therefore \emph{uniquely}
discriminating: a measured propagation delay at the Lieb-Robinson
speed would confirm the physical, finite-velocity spread of
actualization.

\emph{Experimental access.}
Ultracold atoms in optical lattices can be initialized in
controlled superpositions and the spread of decoherence mapped
in real time~\cite{Cheneau2012}. Recent experiments with
superconducting qubit chains~\cite{Mi2022} have demonstrated
Lieb-Robinson light-cone dynamics. Extracting a decoherence
propagation speed from such data would directly test~\eqref{eq:LR}.

\subsection{Prediction 2: Environment-dependent transition duration}

The actualization timescale is set by how long it takes the
environmental entanglement to become irreversible, which is
controlled by the environmental coupling strength.

\textbf{Prediction.}
The duration of a quantum transition (collapse timescale)
scales as
\begin{equation}
  \tauact \sim \frac{\hbar}{\Delta E_{\mathrm{env}}},
  \label{eq:timescale}
\end{equation}
where $\Delta E_{\mathrm{env}}$ is the characteristic energy
scale of the environmental interaction driving the entanglement.
This is not a property of the system alone; it is governed by
the environment.

\emph{Experimental support.}
Minev et al.~\cite{Minev2019} demonstrated that quantum jumps in
superconducting artificial atoms are continuous and governed by
the coupling to the measurement apparatus, consistent
with~\eqref{eq:timescale}.
Pazourek et al.~\cite{Pazourek2015} showed that
photoemission durations depend on the local Coulomb environment.

\emph{The discriminating test.}
Systematically varying the environmental coupling $g$ in a
controlled superconducting circuit---equivalently, tuning
$\Delta E_{\mathrm{env}} \propto g$---should produce a
proportional variation in $\tauact$.
This is a quantitative, parameter-free prediction: the slope
$d\tauact / d(1/\Delta E_{\mathrm{env}}) = \hbar$ is a universal
constant. No free parameters are introduced.

\subsection{Prediction 3: Sharp reversibility threshold}

The criterion~\eqref{eq:criterion} distinguishes two regimes
separated by a physically sharp boundary:
\begin{itemize}
\item \emph{Pre-actualization:} $\Delta S < \Delta S^*$ (threshold
  not yet crossed). The transition is reversible in principle;
  the environmental record is not yet permanent.
\item \emph{Post-actualization:} $\Delta S > \Delta S^*$. The causal
  record is permanent. The transition cannot be reversed regardless
  of any intervention.
\end{itemize}

\textbf{Prediction.}
There exists a sharp threshold $\Delta S^* > 0$ at which
reversibility becomes impossible. Below it, quantum erasure
succeeds; above it, no erasure protocol can restore the
superposition.

\emph{Experimental support.}
Quantum eraser experiments~\cite{ScullyDruhl1982} demonstrate
that which-path information can be erased before the entanglement
becomes irreversible.
Minev et al.~\cite{Minev2019} showed that quantum jumps can be
reversed if the intervention occurs before the jump completes.

\emph{The discriminating test.}
Dynamical collapse models (GRW, CSL) posit a stochastic
collapse with no sharp threshold---the collapse probability
is non-zero at all times, with no clean boundary.
Everettian accounts deny the threshold exists.
The present criterion predicts a \emph{specific, measurable} point
of no return: the entropy production $\Delta S$ crossing a critical
value determined by the interaction Hamiltonian and the vacuum
state geometry.
Mapping this threshold in a superconducting qubit system or
optical cavity would constitute a direct measurement of the
actualization criterion itself.

\begin{table*}[t]
\centering
\caption{Discriminating predictions of the entropic actualization
  criterion against four competing frameworks.
  $\checkmark$ = prediction made; $\times$ = prediction absent or
  contradicted; $\sim$ = partial overlap.
  RA = Relational Actualism (this work); GRW/CSL = dynamical collapse;
  EQM = Everettian many-worlds; RQM = relational QM; Deco.\ = decoherence-only.}
\label{tab:predictions}
\small
\begin{tabular}{lcccccc}
\toprule
Prediction & RA & GRW/CSL & EQM & RQM & Deco.\ \\
\midrule
Finite LR propagation delay ($\Delta t \geq |r_A-r_B|/v_\mathrm{LR}$) & $\checkmark$ & $\times$ & $\times$ & $\times$ & $\times$ \\
Environment-dependent duration ($\tau_\mathrm{act} \sim \hbar/\Delta E_\mathrm{env}$) & $\checkmark$ & $\sim$ & $\times$ & $\times$ & $\sim$ \\
Sharp reversibility threshold ($\Delta S^*$ measurable) & $\checkmark$ & $\times$ & $\times$ & $\times$ & $\times$ \\
\bottomrule
\end{tabular}
\end{table*}

% ── V. DISCUSSION ────────────────────────────────────────────────────────
\section{Discussion}
\label{sec:discussion}

\subsection{Positioning against competing frameworks}

\emph{vs.\ GRW/CSL.}
Dynamical collapse models~\cite{GRW1986,Pearle1989} modify the
Schr\"odinger equation by adding stochastic terms.
They introduce two free parameters ($\lambda$, $r_C$) not
determined by any deeper principle.
The entropic criterion introduces no free parameters: the
threshold is determined entirely by the relative entropy
with respect to the vacuum state.
Furthermore, GRW/CSL predict no Lieb-Robinson propagation delay
and no sharp threshold in the entropic sense.

\emph{vs.\ Everett.}
Everettian QM~\cite{Everett1957} denies that collapse occurs.
All branches co-exist; measurement reveals pre-existing
correlations.
The criterion~\eqref{eq:criterion} identifies a physical
process that distinguishes one branch as actual: the irreversible
entropy production.
The Everettian cannot explain why $\Delta S > 0$ picks out
one branch, because the Everettian framework denies that
any selection occurs.
The three predictions in Section~\ref{sec:predictions}
are therefore directly falsifying for Everett.

\emph{vs.\ Relational QM.}
Rovelli's relational QM~\cite{Rovelli1996} makes collapse
relative to observers.
The criterion~\eqref{eq:criterion} is observer-independent:
the relative entropy with respect to the vacuum $\sigma_0$
is a physical quantity, not relative to any observer.
The Unruh result (Section~\ref{sec:resolution}) demonstrates
this: both inertial and Rindler observers agree on the
$\Delta S = 0$ result.

\emph{vs.\ decoherence-only.}
Decoherence correctly identifies environmental entanglement
as the mechanism for apparent classicality.
The criterion adds the crucial step: it specifies
\emph{which} environmental entanglements constitute
genuine actualizations (those for which $\Delta S > 0$
is \emph{irreversible}).
The sharp reversibility threshold (Prediction~3) is the key
discriminator: decoherence accounts predict no such threshold,
because they deny that genuine collapse occurs.

\subsection{Ontic randomness}

The framework implies irreducibly ontic randomness.
In a two-state superposition $\alpha\ket{x_1} + \beta\ket{x_2}$,
the Born rule probabilities $|\alpha|^2$, $|\beta|^2$ are the
prior distribution over possible actualization locations.
When the actualization occurs at $x_1$ rather than $x_2$, this
is not determined by any prior fact about the system's history.
The actualization event is the unique, ontologically primitive
first link of a new causal chain.
This is not a failure of the framework to explain randomness; it
is the framework's positive claim that the universe is genuinely open.

\subsection{Lean 4 formal verification}

The core finite-dimensional algebraic structure of the criterion
has been formally verified in Lean~4/Mathlib.
Key results include: (i) frame independence---$S(U\rho U^\dagger\|\sigma_0)
= S(\rho\|\sigma_0)$ for any unitary $U$ (theorem
\texttt{frame\_independence} in \texttt{RA\_AQFT\_Proofs\_v10.lean});
(ii) Rindler stationarity---the Rindler thermal state satisfies
$\Delta S = 0$ under modular flow (theorem
\texttt{rindler\_stationarity}); and (iii) the data-processing
inequality (monotonicity of relative entropy under CPTP maps).
The proof files are available at
\href{https://github.com/jsandeman/Relational-Actualism}{github.com/jsandeman/Relational-Actualism}.

% ── VI. CONCLUSION ──────────────────────────────────────────────────────
\section{Conclusion}
\label{sec:conclusion}

We have presented an observer-independent criterion for wave function
collapse: the irreversible increase in quantum relative entropy
$\Delta S(\rho\|\sigma_0) > 0$ with respect to the vacuum state.
The criterion is basis-independent, non-perturbative, and introduces
no free parameters.
It dissolves the measurement problem by identifying collapse with
a physical event (irreversible entropy production) rather than an
observer-relative postulate.
Three experimental predictions discriminate this framework from all
current alternatives: the finite Lieb-Robinson propagation velocity
of actualization, the environment-dependent transition duration
scaling as $\hbar/\Delta E_{\mathrm{env}}$, and the sharp
reversibility threshold $\Delta S^*$.
The Rindler thermal bath is consistent with the criterion:
$\Delta S = 0$ under modular flow, formally verified in Lean~4.

The framework is part of a broader programme---Relational
Actualism---in which the same irreversibility criterion yields
the structural zero of the cosmological constant, a mechanism
for the Hubble tension, and the origin of biological complexity
from a single ontological commitment~\cite{Sandeman2026P3,Sandeman2026P4}.
These consequences are developed in companion papers.

% ── ACKNOWLEDGEMENTS ────────────────────────────────────────────────────
\section*{Acknowledgements}
The author acknowledges the use of Anthropic's Claude and Google's Gemini
as AI research partners during the preparation of this manuscript.
The author is solely responsible for all scientific content and conclusions.

% ── REFERENCES ──────────────────────────────────────────────────────────
\begin{thebibliography}{99}

\bibitem{Bell1990}
J.~S.~Bell,
\emph{Against ``Measurement''},
Phys.\ World \textbf{3}, 33 (1990).

\bibitem{GRW1986}
G.~C.~Ghirardi, A.~Rimini, and T.~Weber,
Unified dynamics for microscopic and macroscopic systems,
Phys.\ Rev.\ D \textbf{34}, 470 (1986).

\bibitem{Pearle1989}
P.~Pearle,
Combining stochastic dynamical state-vector reduction with spontaneous
localization,
Phys.\ Rev.\ A \textbf{39}, 2277 (1989).

\bibitem{Everett1957}
H.~Everett III,
``Relative state'' formulation of quantum mechanics,
Rev.\ Mod.\ Phys.\ \textbf{29}, 454 (1957).

\bibitem{Rovelli1996}
C.~Rovelli,
Relational quantum mechanics,
Int.\ J.\ Theor.\ Phys.\ \textbf{35}, 1637 (1996).

\bibitem{Bohm1952}
D.~Bohm,
A suggested interpretation of the quantum theory in terms of ``hidden''
variables I and II,
Phys.\ Rev.\ \textbf{85}, 166, 180 (1952).

\bibitem{Joos1985}
E.~Joos and H.~D.~Zeh,
The emergence of classical properties through interaction with the
environment,
Z.\ Phys.\ B \textbf{59}, 223 (1985).

\bibitem{Zurek2003}
W.~H.~Zurek,
Decoherence, einselection, and the quantum origins of the classical,
Rev.\ Mod.\ Phys.\ \textbf{75}, 715 (2003).

\bibitem{Unruh1976}
W.~G.~Unruh,
Notes on black-hole evaporation,
Phys.\ Rev.\ D \textbf{14}, 870 (1976).

\bibitem{Umegaki1962}
H.~Umegaki,
Conditional expectation in an operator algebra IV (entropy and information),
K\=odai Math.\ Sem.\ Rep.\ \textbf{14}, 59 (1962).

\bibitem{Lindblad1975}
G.~Lindblad,
Completely positive maps and entropy inequalities,
Commun.\ Math.\ Phys.\ \textbf{40}, 147 (1975).

\bibitem{Haag1996}
R.~Haag,
\emph{Local Quantum Physics}, 2nd ed.\
(Springer, Berlin, 1996).

\bibitem{LiebRobinson1972}
E.~H.~Lieb and D.~W.~Robinson,
The finite group velocity of quantum spin systems,
Commun.\ Math.\ Phys.\ \textbf{28}, 251 (1972).

\bibitem{Cheneau2012}
M.~Cheneau et al.,
Light-cone-like spreading of correlations in a quantum many-body system,
Nature \textbf{481}, 484 (2012).

\bibitem{Mi2022}
X.~Mi et al.,
Information scrambling in quantum circuits,
Science \textbf{374}, 1479 (2021).

\bibitem{Minev2019}
Z.~K.~Minev et al.,
To catch and reverse a quantum jump mid-flight,
Nature \textbf{570}, 200 (2019).

\bibitem{Pazourek2015}
R.~Pazourek, S.~Nagele, and J.~Burgd\"{o}rfer,
Attosecond chronoscopy of photoemission,
Rev.\ Mod.\ Phys.\ \textbf{87}, 765 (2015).

\bibitem{ScullyDruhl1982}
M.~O.~Scully and K.~Dr\"{u}hl,
Quantum eraser: A proposed photon correlation experiment concerning
observation and ``delayed choice'' in quantum mechanics,
Phys.\ Rev.\ A \textbf{25}, 2208 (1982).

\bibitem{Sandeman2026P3}
J.~F.~Sandeman,
Structural Zero Cosmological Constant and the Hubble Tension from
Causal Actualization Density,
in preparation (2026).

\bibitem{Sandeman2026P4}
J.~F.~Sandeman,
Relational Actualism: Irreversible Events as the Primitive Basis of
Physical Reality,
in preparation (2026).

\end{thebibliography}

\end{document}
